\FloatBarrier
\section{Limitations}
\label{app:limitations}

% 中文：本文主要研究同一物理场景中由视角变化引起的镜头切换。面对复杂剪辑、人物反复进出画面、
% 快速画面变化以及很长的视频，持续保持空间与身份一致性仍然具有挑战性。
Our evaluation focuses on shot transitions caused by viewpoint changes within
the same physical scene. Maintaining spatial and identity consistency under
complex editing, repeated entrances and exits, rapid visual changes, and very
long videos remains challenging.

% 中文：流式与离线方法的时间信息不同，因此本文分别说明各方法是否使用未来帧。
% 此外，条件人体误差需要与 Coverage 一同阅读；更好的相机或身份一致性也不保证每一个人体指标都改善。
Streaming and offline methods have different temporal access, which is stated
throughout the comparisons. Human errors computed only on valid matches should
also be read together with Coverage. Improved camera or identity consistency
does not imply improvement in every body metric; the complete Harmony4D
results retain these metric-level differences.
